\section{决策概率与治理建议}

有序Logistic回归提供的系数与OR值能够刻画影响方向与强度，但对治理决策而言，决策概率与边际效应更为直观。为将统计结果转化为可操作的治理依据，本节基于有序Logistic模型计算决策概率与边际效应。

\subsection{基线决策概率}

基于有序Logistic模型，在全样本特征均值处计算得到的基线决策概率为：居民对声环境感到“不满意”（$Y\le2$）的概率为44.84\%，“满意”（$Y\ge4$）的概率为25.34\%，期望满意度 $E[Y]=2.78$。这表明整体而言居民对声环境的评价集中于“一般”（$Y=3$）附近，不满意人群占比接近一半，治理空间较大。

\subsection{可干预变量的平均边际效应}

为识别治理的优先级，本文计算了各可干预变量对“不满意”概率 $P(Y\le2)$ 的平均边际效应（Average Marginal Effect，AME），结果如表\ref{tab:marginal}所示。

\begin{table}[H]
\centering
\caption{各可干预变量对不满意概率的平均边际效应}
\label{tab:marginal}
\rowcolors{2}{tablight}{white}
\begin{tabular}{M{7cm}M{4cm}}
\bluetoprule
\rowcolor{tabhead}
\thead{干预变量} & \thead{$\Delta P(Y\le2)$（百分点）} \\
\hline
隔音墙（无 $\to$ 有） & $-13.76$ \\
县城（vs 乡镇农村） & $+8.20$ \\
郊区（vs 乡镇农村） & $+4.46$ \\
杭州（vs 温州） & $+3.33$ \\
主城区（vs 乡镇农村） & $-3.92$ \\
宁波（vs 温州） & $-1.68$ \\
\bluebottomrule
\end{tabular}
\end{table}

由表\ref{tab:marginal}可知，隔音墙的边际效应最强：住宅配备隔音墙可使居民“不满意”概率平均下降13.76个百分点，是当前最具杠杆效应的干预手段。杭州相对于温州“不满意”概率高3.33个百分点，是三市中治理压力较大的城市。需要说明的是，主城区相对于乡镇农村的边际效应为负（$-3.92$个百分点），方向与直觉相反，系数据缺陷同源所致，本文不对其作过度解读。

\subsection{群体决策概率差异}

进一步按关键群体分列决策概率，结果如表\ref{tab:groupprob}所示。

\begin{table}[H]
\centering
\caption{关键群体的决策概率}
\label{tab:groupprob}
\rowcolors{2}{tablight}{white}
\begin{tabular}{M{3.2cm}M{3.8cm}M{3.8cm}M{3.4cm}}
\bluetoprule
\rowcolor{tabhead}
\thead{分组} & \thead{类别} & \thead{$P(Y\le2)$ 不满意概率} & \thead{$P(Y\ge4)$ 满意概率} \\
\hline
\multirow{3}{*}{城市} & 杭州 & 49.61\% & 21.40\% \\
 & 宁波 & 42.47\% & 27.97\% \\
 & 温州 & 43.64\% & 25.45\% \\
\hline
\multirow{2}{*}{隔音墙} & 无/不清楚 & 50.24\% & 20.61\% \\
 & 有 & 31.15\% & 37.34\% \\
\hline
\multirow{5}{*}{年龄} & 18岁以下 & 58.68\% & 15.53\% \\
 & 19-30岁 & 35.85\% & 33.63\% \\
 & 31-45岁 & 45.28\% & 24.70\% \\
 & 46-60岁 & 50.11\% & 20.92\% \\
 & 61岁及以上 & 44.60\% & 22.80\% \\
\bluebottomrule
\end{tabular}
\end{table}

由表\ref{tab:groupprob}可知，18岁以下群体的“不满意”概率最高（58.68\%），显著高于其他年龄组；无隔音墙群体（50.24\%）的不满意概率明显高于有隔音墙群体（31.15\%）；杭州居民的不满意概率（49.61\%）为三市最高。

\subsection{治理建议的量化依据}

综合边际效应与群体概率分析，可凝练出以下治理优先级：隔音改造是第一优先级工程手段，其边际效应远高于其他变量；青少年（18岁以下）是重点人群，其不满意概率最高；杭州是重点治理城市。上述量化结论为后文的优化建议提供了实证支撑。

